\clearpage
\fancyhead[L]{\small THEORIA INTERPOLATIONIS}
\pdfbookmark[0]{Theoria interpolationis methodo nova tractata}{interpolation-nachlass}
\begingroup
\renewcommand{\arraystretch}{1.35}
\SourceLeaf{274}
\begin{center}\textbf{NACHLASS.}\par\bigskip
\textbf{THEORIA INTERPOLATIONIS}\par\medskip
\textbf{METHODO NOVA TRACTATA.}\par\medskip
\rule{.15\linewidth}{.4pt}\end{center}
\InterpolationArticle{1}
\textsc{Problema.}\quad\textit{Invenire summam seriei}
\[
\begin{aligned}
&\frac{a^n}{(a-b)(a-c)(a-d)(a-e)\ldots}\\[5pt]
&+\frac{b^n}{(b-a)(b-c)(b-d)(b-e)\ldots}\\
&+\frac{c^n}{(c-a)(c-b)(c-d)(c-e)\ldots}\\[5pt]
&+\frac{d^n}{(d-a)(d-b)(d-c)(d-e)\ldots}\\
&+\frac{e^n}{(e-a)(e-b)(e-c)(e-d)\ldots}+\text{etc.}
\end{aligned}
\]
\textit{ubi $a,b,c,d,e$ sunt $m$ quantitates diversae, atque $n$ numerus integer quicunque positivus, negativus sive etiam 0.}

\textit{Solutio.}\quad Faciendo brevitatis caussa
\[
\begin{aligned}
\frac1{(a-b)(a-c)(a-d)(a-e)\ldots}&=\alpha\\
\frac1{(b-a)(b-c)(b-d)(b-e)\ldots}&=\beta\\
\frac1{(c-a)(c-b)(c-d)(c-e)\ldots}&=\gamma\\
\frac1{(d-a)(d-b)(d-c)(d-e)\ldots}&=\delta,\ \text{etc.}
\end{aligned}
\]
ita ut summa quaesita, quam per $S^n$ denotabimus fiat $=\alpha a^n+\beta b^n+\gamma c^n+\delta d^n+$ etc.: manifestum est, si $x$ exprimat quantitatem indeterminatam, ex evolutione aggregati
\SourceLeaf{275}
\[
P=\frac\alpha{1-ax}+\frac\beta{1-bx}+\frac\gamma{1-cx}+\frac\delta{1-dx}+\text{etc.}
\]
in seriem secundum potestates ipsius $x$ ascendentem, prodire
\[
S^0+S^1x+S^2xx+S^3x^3+\text{etc.}\quad\text{in infn.}
\]
Statuatur $(1-ax)(1-bx)(1-cx)(1-dx)\ldots=Q$, eritque $Q$ functio integra indeterminatae $x$, ad ordinem $m^{\mathrm{tum}}$ ascendens: $PQ$ autem fiet functio integra ordinis $m-1^{\mathrm{ti}}$ puta $=$
\[
\begin{aligned}
&\alpha(1-bx)(1-cx)(1-dx)\ldots\\
&+\beta(1-ax)(1-cx)(1-dx)\ldots\\
&+\gamma(1-ax)(1-bx)(1-dx)\ldots\\
&+\delta(1-ax)(1-bx)(1-cx)\ldots\\
&+\text{etc.}
\end{aligned}
\]
Qua propius considerata, patebit, per substitutionem $x=\frac1a$ omnes partes praeter primam evanescere, hanc vero abire in
\[
\alpha(1-\tfrac ba)(1-\tfrac ca)(1-\tfrac da)\ldots=\frac1{a^{m-1}}
\]
Simili modo per substitutionem $x=\frac1b$ evanescent omnes partes praeter secundam, quae fit $=\frac1{b^{m-1}}$. Perinde per substitutiones $x=\frac1c$, $x=\frac1d$ etc. transit $PQ$ in $\frac1{c^{m-1}}$, $\frac1{d^{m-1}}$ etc. Hinc vero sequitur, $PQ-x^{m-1}$ per omnes has substitutiones valorem 0 obtinere, quod fieri nequit, nisi fuerit identice $=0$, sive $PQ=x^{m-1}$; alioquin enim aequatio $PQ-x^{m-1}=0$, quae non maioris quam $m-1^{\mathrm{ti}}$ ordinis est, $m$ radices diversas $\frac1a,\frac1b,\frac1c,\frac1d\ldots$ haberet.

Iam sit
\[
Q=1-Ax+Bxx-Cx^3+Dx^4-\text{etc.}
\]
nempe $A$ summa quantitatum $a,b,c,d\ldots$; $B$ summa productorum e binis; $C$ summa productorum e ternis etc., patetque, quum ex evolutione fractionis $\frac{x^{m-1}}Q=P$ prodire debeat
\[
S^0+S^1x+S^2xx+S^3x^3+\text{etc.}
\]
primo: esse debere $S^0=0$, $S^1=0$, $S^2=0$ etc. usque ad $S^{m-2}=0$\HistoricalNote{22}, tunc
\SourceLeaf{276}
vero fieri $S^{m-1}=1$, $S^m=A$, tandemque terminos ulteriores tamquam membra seriei recurrentis per legem sequentem determinari:
\[
\begin{aligned}
S^m&=A\\
S^{m+1}&=AS^m-BS^{m-1}=AA-B\\
S^{m+2}&=AS^{m+1}-BS^m+CS^{m-1}=A^3-2BA+C\\
&\text{etc.}
\end{aligned}
\]
Facile quidem hinc colligitur $S^{m+1}$ esse $=aa+bb+cc+\ldots+B$ sive summam quadratorum cum summa omnium productorum e binis diversis quantitatum $a,b,c,d\ldots$; sed quo clarius perspiciatur, quonam modo termini sequentes ex elementis $a,b,c,d\ldots$ formentur, observamus $\frac1Q$ esse productum e seriebus
\[
\begin{aligned}
&1+ax+aaxx+a^3x^3+\text{etc.}\\
&1+bx+bbxx+b^3x^3+\text{etc.}\\
&1+cx+ccxx+c^3x^3+\text{etc.}\\
&1+dx+ddxx+d^3x^3+\text{etc.}\\
&\text{etc.}
\end{aligned}
\]
Hoc vero productum est $=\Sigma a^\lambda b^\mu c^\nu\ldots\times x^{\lambda+\mu+\nu.\,.}$, ubi exponentibus $\lambda,\mu,\nu\ldots$ omnes valores integri a 0 usque in infn. tribuendi omnibusque quibus fieri potest modis combinandi sunt. Quocirca ut in serie, in quam $\frac1Q$ evolvitur, eius termini, qui continet $x^{n+1-m}$, coëfficientem obtineamus, numerum $n+1-m$ omnibus quibus fieri potest modis in $n$ partes integras $\lambda+\mu+\nu+.\,.$ (inter quas etiam pars 0 admittitur) discerpere oportet, omnibus quoque permutationibus harum partium permissis; atque tunc singula producta $a^\lambda b^\mu c^\nu\ldots$ in summam colligere, quae erit coëfficiens quaesitus, simulque $=S^n$. Levi attentione adhibita patebit, huic regulae prorsus aequivalere sequentem: Ex $m$ quantitatibus $a,b,c,d\ldots$ omnes combinationes $n+1-m$ elementorum colligendae, admissis repetitionibus, et singulae tamquam producta considerandae, quorum aggregatum erit $=S^n$. Quare erit ut supra $S^{m+1}$ summa omnium productorum e binis quantitatum $a,b,c,d\ldots$ tum diversis tum identicis; $S^{m+2}$ summa omnium productorum e ternis diversis seu identicis etc.

Nihil iam superest, nisi ut summam progressionis nostrae pro valoribus negativis ipsius $n$ definire doceamus. Ad quem finem partem primam summae $S^{-n}$, puta $\frac{a^{-n}}{(a-b)(a-c)(a-d)\ldots}$ sub hanc formam ponemus
\SourceLeaf{277}
\[
\pm\frac1{abcd\ldots}\times\frac{(\frac1a)^{m+n-2}}{(\frac1a-\frac1b)(\frac1a-\frac1c)(\frac1a-\frac1d)\ldots}
\]
ubi signum superius vel inferius adoptandum est, prout $m$ impar est vel par, similisque transformatio etiam ad partes reliquas applicari poterit. Quamobrem si per characterem $T$ designetur expressio, quae perinde ex $\frac1a,\frac1b,\frac1c,\frac1d.\,.$ oritur, ut $S$ ex $a,b,c,d\ldots$, manifesto fiet $S^{-n}=\pm\frac{T^{m+n-2}}{abcd\ldots}$. Hoc itaque modo hic casus ad praecedentem reductus est, fitque
\[
S^{-1}=\pm\frac1{abcd\ldots},\qquad S^{-2}=\pm\frac{\frac1a+\frac1b+\frac1c+\frac1d+.\,.}{abcd\ldots},
\]
$S^{-3}$ aequalis producto ex $\pm\frac1{abcd\ldots}$ in summam omnium productorum e binis quantitatum $\frac1a,\frac1b,\frac1c,\frac1d\ldots$ diversis aut identicis etc.
\InterpolationArticle{2}
Applicabimus disquisitionem praecedentem ad eum casum, ubi quantitatibus $a,b,c,d$ valores imaginarii tribuuntur: hac ratione ad quasdam insignes relationes perfacile perveniemus, quae alia methodo tractatae maiores difficultates obiicerent. Sit $E$ basis logarithmorum naturalium, $i$ quantitas imaginaria $\BareRoot{-1}$; consideremus loco quantitatum realium $a,b,c\ldots$ imaginarias $E^{ia},E^{ib},E^{ic},E^{id}\ldots$ et $E^{-ia},E^{-ib},E^{-ic},E^{-id}\ldots$, statuamusque
\[
\begin{aligned}
&\frac{E^{ina}}{(E^{ia}-E^{ib})(E^{ia}-E^{ic})(E^{ia}-E^{id})\ldots}\\
&+\frac{E^{inb}}{(E^{ib}-E^{ia})(E^{ib}-E^{ic})(E^{ib}-E^{id})\ldots}\\
&+\frac{E^{inc}}{(E^{ic}-E^{ia})(E^{ic}-E^{ib})(E^{ic}-E^{id})\ldots}\\
&+\frac{E^{ind}}{(E^{id}-E^{ia})(E^{id}-E^{ib})(E^{id}-E^{ic})\ldots}+\text{etc.}=S^n,\quad\text{atque}
\end{aligned}
\]
\[
\begin{aligned}
&\frac{E^{-ina}}{(E^{-ia}-E^{-ib})(E^{-ia}-E^{-ic})(E^{-ia}-E^{-id})\ldots}\\
&+\frac{E^{-inb}}{(E^{-ib}-E^{-ia})(E^{-ib}-E^{-ic})(E^{-ib}-E^{-id})\ldots}\\
&+\frac{E^{-inc}}{(E^{-ic}-E^{-ia})(E^{-ic}-E^{-ib})(E^{-ic}-E^{-id})\ldots}\\
&+\frac{E^{-ind}}{(E^{-id}-E^{-ia})(E^{-id}-E^{-ib})(E^{-id}-E^{-ic})\ldots}+\text{etc.}=T^n
\end{aligned}
\]
\SourceLeaf{278}
Designando itaque multitudinem quantitatum $a,b,c,d\ldots$ per $m$, erunt $S^0,S^1,S^2\ldots S^{m-2}$, nec non $T^0,T^1,T^2\ldots T^{m-2}$ omnes $=0$; porro $S^{m-1}=T^{m-1}=1$; $S^m$ summa quantitatum $E^{ia},E^{ib},E^{ic},E^{id}\ldots$; $S^{m+1}$ summa productorum omnium e binis diversis seu identicis; $S^{m+2}$ summa productorum e ternis etc.; et perinde valores summarum $T^m,T^{m+1},T^{m+2}$ etc. e quantitatibus $E^{-ia},E^{-ib},E^{-ic},E^{-id}\ldots$ formandi erunt.

Iam quum constet, esse $E^{ix}+E^{-ix}=2\cos x$, $E^{ix}-E^{-ix}=2i\sin x$, facile perspicietur, valorem expressionis $\frac12(S^n+T^n)$, qui pro $n=0,1,2\ldots m-2$ fit $=0$, pro $n=m-1$ autem $=1$, pro $n=m$ fieri
\[
=\cos a+\cos b+\cos c+\cos d\,.\,.\,.\,.
\]
similiterque fieri $\frac12(S^{m+1}+T^{m+1})$ summam cosinuum omnium angulorum, qui oriuntur addendo \textit{binos} ex his $a,b,c,d.\,.$ diversos seu identicos; $\frac12(S^{m+2}+T^{m+2})$ summam cosinuum omnium angulorum, qui oriuntur addendo ex iisdem \textit{ternos} etc. Perinde erit $\frac{S^n-T^n}{2i}=0$ pro $n=0,1,2\ldots m-1$; porro $=\sin a+\sin b+\sin c+\sin d+.\,.$ pro $n=m$; et similiter $\frac{S^{m+1}-T^{m+1}}{2i}$ erit summa sinuum omnium angulorum, qui oriuntur addendo ex his $a,b,c,d\ldots$ binos diversos seu identicos; $\frac{S^{m+2}-T^{m+2}}{2i}$ summa sinuum omnium angulorum, qui oriuntur ex iisdem, ternos combinando etc.

Summarum $S^n,T^n$ partes nunc propius considerabimus. Est
\[
\begin{aligned}E^{ia}-E^{ib}&=E^{\frac12i(a+b)}(E^{\frac12i(a-b)}-E^{\frac12i(b-a)})\\
&=2iE^{\frac12i(a+b)}\sin\tfrac12(a-b)\end{aligned}
\]
perinde
\[
E^{ia}-E^{ic}=2iE^{\frac12i(a+c)}\sin\tfrac12(a-c)\quad\text{etc.}
\]
Quamobrem in $S^n$ partis primae denominator fit
\[
=(2i)^{m-1}E^{\frac12i(b+c+d\ldots+(m-1)a)}\sin\tfrac12(a-b)\sin\tfrac12(a-c)\sin\tfrac12(a-d)\ldots
\]
statuendoque $a+b+c+d+\ldots=s$, haec pars ipsa
\[
=\frac{E^{i((n+1-\frac12m)a-\frac12s)}}{(2i)^{m-1}\sin\tfrac12(a-b)\sin\tfrac12(a-c)\sin\tfrac12(a-d)\ldots}
\]
Simili modo habetur
\[
E^{-ia}-E^{-ib}=-2iE^{-\frac12i(a+b)}\sin\tfrac12(a-b)
\]
unde tandem pars prima summae $T^n$ provenit
\SourceLeaf{279}
\[
=\frac{\pm E^{-i((n+1-\frac12m)a-\frac12s)}}{(2i)^{m-1}\sin\tfrac12(a-b)\sin\tfrac12(a-c)\sin\tfrac12(a-d)\ldots}
\]
ubi signum superius vel inferius valet, prout $m$ impar est vel par. Quam cum parte prima summae $S^n$ addendo, sive ab eadem subtrahendo, concludimus fieri \textit{primo} pro valore impari ipsius $m$ partem primam summae $\frac12(S^n+T^n)$
\[
=\frac{\cos((n+1-\frac12m)a-\frac12s)}{(2i)^{m-1}\sin\tfrac12(a-b)\sin\tfrac12(a-c)\sin\tfrac12(a-d)\ldots}
\]
partem primam summae $\frac{S^n-T^n}{2i}$
\[
=\frac{\sin((n+1-\frac12m)a-\frac12s)}{(2i)^{m-1}\sin\tfrac12(a-b)\sin\tfrac12(a-c)\sin\tfrac12(a-d)\ldots}
\]
\textit{secundo} pro valore pari ipsius $m$ partem primam summae $\frac12(S^n+T^n)$
\[
=\frac{\sin((n+1-\frac12m)a-\frac12s)}{2^{m-1}i^{m-2}\sin\tfrac12(a-b)\sin\tfrac12(a-c)\sin\tfrac12(a-d)\ldots}
\]
partem primam summae $\frac{S^n-T^n}{2i}$
\[
=\frac{-\cos((n+1-\frac12m)a-\frac12s)}{2^{m-1}i^{m-2}\sin\tfrac12(a-b)\sin\tfrac12(a-c)\sin\tfrac12(a-d)\ldots}
\]
Manifesto partes sequentes expressionum $\frac12(S^n+T^n)$, $\frac{S^n-T^n}{2i}$ primae prorsus analogae erunt, atque inde per solam commutationem characteris $a$ cum $b,c,d\ldots$ orientur. Iam designando per $k$ angulum arbitrarium ponendoque $n+1-\frac12m=\lambda$, adiumento formularum
\[
\begin{aligned}
\cos(\lambda a+k)&=\cos(k+\tfrac12s)\cos(\lambda a-\tfrac12s)\\&\qquad-\sin(k+\tfrac12s)\sin(\lambda a-\tfrac12s)\\
\sin(\lambda a+k)&=\cos(k+\tfrac12s)\sin(\lambda a-\tfrac12s)\\&\qquad+\sin(k+\tfrac12s)\cos(\lambda a-\tfrac12s)
\end{aligned}
\]
haud difficile perveniemus ad summationem serierum sequentium
\[
\begin{aligned}
&\frac{\cos(\lambda a+k)}{\sin\tfrac12(a-b)\sin\tfrac12(a-c)\sin\tfrac12(a-d)\ldots}\\
&+\frac{\cos(\lambda b+k)}{\sin\tfrac12(b-a)\sin\tfrac12(b-c)\sin\tfrac12(b-d)\ldots}\\
&+\frac{\cos(\lambda c+k)}{\sin\tfrac12(c-a)\sin\tfrac12(c-b)\sin\tfrac12(c-d)\ldots}\\
&+\frac{\cos(\lambda d+k)}{\sin\tfrac12(d-a)\sin\tfrac12(d-b)\sin\tfrac12(d-c)\ldots}+\text{etc.}=U^\lambda
\end{aligned}
\]
atque
\SourceLeaf{280}
\[
\begin{aligned}
&\frac{\sin(\lambda a+k)}{\sin\tfrac12(a-b)\sin\tfrac12(a-c)\sin\tfrac12(a-d)\ldots}\\
&+\frac{\sin(\lambda b+k)}{\sin\tfrac12(b-a)\sin\tfrac12(b-c)\sin\tfrac12(b-d)\ldots}\\
&+\frac{\sin(\lambda c+k)}{\sin\tfrac12(c-a)\sin\tfrac12(c-b)\sin\tfrac12(c-d)\ldots}\\
&+\frac{\sin(\lambda d+k)}{\sin\tfrac12(d-a)\sin\tfrac12(d-b)\sin\tfrac12(d-c)\ldots}+\text{etc.}=V^\lambda.
\end{aligned}
\]
\begin{center}\textit{Casus primus, si $m$ est impar.}\end{center}
In hoc casu $\lambda$ debet esse non numerus integer sed fractione $\frac12$ affectus sive numeri imparis semissis. Ad valores \textit{negativos} ipsius $\lambda$ non opus est respicere, quum habeatur $\cos(-\lambda a+k)=\cos(\lambda a-k)$ atque $\sin(-\lambda a+k)=-\sin(\lambda a-k)$, adeoque summatio pro valore negativo e summatione pro opposito positivo per solam mutationem ipsius $k$ in $-k$ sponte demanat: haec observatio manifesto etiam in casu sequenti valebit. Iam patet facile, pro $n=0,1,\ldots m-2$, sive pro $\lambda=-\frac12m+1,-\frac12m+2,\ldots\frac12m-1$, sive ut valores negativos omittamus, pro $\lambda=\frac12,\frac32,\frac52\ldots\frac12m-1$ fieri tum $U^\lambda=0$, tum $V^\lambda=0$; porro
\[
\begin{aligned}
U^{\frac12m}&=(2i)^{m-1}\cos(\tfrac12s+k),\\V^{\frac12m}&=(2i)^{m-1}\sin(\tfrac12s+k)\\
U^{\frac12m+1}&=(2i)^{m-1}\{\cos(\tfrac12s+k+a)+\cos(\tfrac12s+k+b)\\&\qquad+\cos(\tfrac12s+k+c)\\
&\qquad+\cos(\tfrac12s+k+d)+\text{etc.}\}\\
V^{\frac12m+1}&=(2i)^{m-1}\{\sin(\tfrac12s+k+a)+\sin(\tfrac12s+k+b)\\&\qquad+\sin(\tfrac12s+k+c)\\
&\qquad+\sin(\tfrac12s+k+d)+\text{etc.}\}
\end{aligned}
\]
$U^{\frac12m+2}$ vel $V^{\frac12m+2}$ aequalem producto ex $(2i)^{m-1}$ in summam cosinuum vel sinuum omnium angulorum, qui oriuntur addendo $\frac12s+k$ cum binis angulorum $a,b,c,d\ldots$ diversis seu identicis; $U^{\frac12m+3}$ vel $V^{\frac12m+3}$ aequalem producto ex $(2i)^{m-1}$ in summam cosinuum vel sinuum omnium angulorum, qui oriuntur addendo $\frac12s+k$ cum ternis ex iisdem etc. Ceterum vix opus erit admonere, potestatem quantitatis imaginariae $i$ cum exponente pariter pari fieri $=+1$, cum impariter pari $=-1$.
\begin{center}\textit{Casus secundus, si $m$ est par.}\end{center}
In hoc casu $\lambda$ debet esse numerus integer, fitque pro $\lambda=0,1,2\ldots\frac12m-1$ tum $U^\lambda=0$, tum $V^\lambda=0$. Porro erit
\SourceLeaf{281}
\[
\begin{aligned}
U^{\frac12m}&=2^{m-1}i^{m-2}\sin(\tfrac12s+k),\\
V^{\frac12m}&=-2^{m-1}i^{m-2}\cos(\tfrac12s+k);\\
U^{\frac12m+1}&=2^{m-1}i^{m-2}\{\sin(\tfrac12s+k+a)+\sin(\tfrac12s+k+b)\\&\qquad+\sin(\tfrac12s+k+c)\\
&\qquad+\sin(\tfrac12s+k+d)+\text{etc.}\}\\
V^{\frac12m+1}&=-2^{m-1}i^{m-2}\{\cos(\tfrac12s+k+a)+\cos(\tfrac12s+k+b)\\&\qquad+\cos(\tfrac12s+k+c)\\
&\qquad+\cos(\tfrac12s+k+d)+\text{etc.}\}
\end{aligned}
\]
$U^{\frac12m+2},U^{\frac12m+3}$ etc. aequalis producto ex $2^{m-1}i^{m-2}$ in summam sinuum omnium angulorum, qui oriuntur addendo $\frac12s+k$ cum binis, ternis etc. angulorum $a,b,c\ldots$ diversis seu identicis; $V^{\frac12m+2},U^{\frac12m+3}$ etc. aequalis producto ex $-2^{m-1}i^{m-2}$ in summam cosinuum omnium angulorum, qui oriuntur addendo $\frac12s+k$ cum binis, ternis etc. eorundem angulorum.
\InterpolationArticle{3}
Sit $X$ functio indeterminatae $x$ huius formae
\[
\alpha+\beta x+\gamma xx+\delta x^3+\text{etc.}
\]
quae non excurrat in infinitum, sed abrumpatur, neque ultra terminum, qui continet $x^{m-1}$ egrediatur, ita ut multitudo coëfficientium non sit maior quam $m$. Tunc si pro $m$ valoribus diversis ipsius $x$, puta $a,b,c,d\ldots$ valores correspondentes ipsius $X$ sunt cogniti, puta $=A,B,C,D\ldots$ ex his valor ipsius $X$ valori alicui alii ipsius $x$ respondens sequenti modo concinne eruetur. Sit $t$ valor novus ipsius $x$, atque $T$ valor respondens ipsius $X$, ita ut sequentes $m+1$ aequationes locum habeant
\[
\begin{aligned}
A&=\alpha+\beta a+\gamma aa+\delta a^3+\ldots\\
B&=\alpha+\beta b+\gamma bb+\delta b^3+\ldots\\
C&=\alpha+\beta c+\gamma cc+\delta c^3+\ldots\\
D&=\alpha+\beta d+\gamma dd+\delta d^3+\ldots\\
&\text{etc.}\\
T&=\alpha+\beta t+\gamma tt+\delta t^3+\ldots
\end{aligned}
\]
\endgroup
